%! TeX root = ./main.tex
% vim: set spelllang=en:
\section{sowing --- Manipulating directed acyclic graphs}
\programdetails{https://github.com/UdeM-LBIT/sowing}{GPL v3.0}{Between May~2023 and April~2026, 1\,400 lines of code}

\emph{sowing} is a library written in the Python language that offers a persistent data structure for efficiently representing, building, traversing, modifying and serializing trees, forests and, more generally, any directed acyclic graph.
The advantages of using such a data structure are manifold:
\begin{compactitemize}
    \item common subgraphs can be reused, thus reducing the memory footprint, without however risking bugs related to the modification of a shared fragment;
    \item graphs can be passed from one function or thread of execution to another without requiring a copy;
    \item graphs can be hashed, which allows them to be used as elements of other data structures, such as dictionaries or sets;
    \item the absence of side effects makes it possible to adopt a more functional programming style, which makes it easier to reason about the resulting programs.
\end{compactitemize}

Graphs are represented in their usual recursive form: the \texttt{Node} class represents a node, including the data attached to it, and the outgoing arcs towards the successors of the node, themselves represented by instances of \texttt{Node}.
In directed acyclic graphs that are not trees, nodes with several predecessors are represented by making several nodes point to the same object in memory.
Forests can be represented with the \texttt{Hedge} class, which attaches to them an artificial node acting as a root.

A \emph{cursor} can be created from any node.
This cursor makes it possible to move freely in the graph while keeping track of the context around the current node.
Thanks to this structure, it is possible to perform local modifications of the tree or graph in constant time, despite the immutable nature of the structure (see \cref{lst:logiciel:sowing-ex} for an example).
This structure is based on the \emph{zipper} principle introduced by \textcite{Huet1997Zipper}.

\begin{listing}[h]
\begin{minted}{python}
def display(root: Node, leaves: set[str]) -> Node:
  def filter_tree(cursor: Zipper) -> Zipper:
    node = cursor.node
    if cursor.is_leaf() and node.data["name"] not in leaves:
      # removal of the targeted leaves
      return cursor.replace(node=None)
    if len(node.edges) == 1:
      # contraction of unary nodes
      return cursor.replace(node=node.edges[0].node)
    return cursor
  return traversal.fold(filter_tree, traversal.depth(root))

tree = newick.parse("((a,b,c),(d,e));")
disp = display(tree, {"a", "b", "d"})
newick.write(disp)
# result: ((a,b),d);
\end{minted}
    \caption[Using the \emph{sowing} library to modify a tree]{
        Function that contracts a tree to keep a given set of leaves.
        Despite the immutable nature of the data structure, this function uses \(\bigo(n)\) time and space, where \(n\) is the number of nodes in the initial tree.
        \label{lst:logiciel:sowing-ex}
    }
\end{listing}

The features of \emph{sowing} are divided among several modules.

\begin{itemize}
    \item \textbf{\texttt{sowing.node}}, \textbf{\texttt{sowing.zipper}}, \textbf{\texttt{sowing.hedge}} contain the main data structure through the classes \texttt{Node}, \texttt{Zipper} and \texttt{Hedge}.
    \item \textbf{\texttt{sowing.indexed}} defines \texttt{IndexedTree}, a structure that wraps a labelled tree to allow constant-time access to any of its nodes by its label.
        This structure also makes it possible to perform the following operations in constant time: check whether two nodes are ancestors of one another, compute the distance between two nodes and identify the lowest common ancestor or \emph{lca} of two nodes.
        This structure is based on a linearithmic-time preprocessing described by \textcite{Bender2000LcaProblemRevisited}.
    \item \textbf{\texttt{sowing.traversal}} provides iterators for traversing a tree or directed acyclic graph structure in different orders: depth-first preorder or postorder, topological order or Euler order.
        The \texttt{fold} and \texttt{map} functions make it possible to modify a tree in linear time following any of these traversal orders.
    \item \textbf{\texttt{sowing.repr}} makes it possible to load or serialize trees and
        directed acyclic graphs in several formats: Newick \parencite{Felsenstein1986NewickTreeFormat}, JSON or DOT (Graphviz).
    \item \textbf{\texttt{sowing.comb.binary}} makes it possible to enumerate the resolutions of a non-binary tree without considering the order of the successors of the different nodes.
    \item \textbf{\texttt{sowing.comb.supertree}} implements the algorithm of \textcite{Ng1996ReconstructionRootedTrees} for building supertrees that extend several phylogenetic trees.
\end{itemize}

